\newcommand{\hmwkTitle}{Homework 0}
\newcommand{\hmwkDueDate}{Tuesday, January 20, 2026}
\newcommand{\hmwkHeaderTitle}{ICCS206 Discrete Mathematic: Homework 0}
\documentclass{article}

\usepackage{fancyhdr}
\usepackage{amsmath}
\allowdisplaybreaks
\usepackage{amsthm}
\usepackage{amsfonts}
\usepackage{tikz}
\usepackage[plain]{algorithm}
\usepackage{algpseudocode}
\usepackage{enumitem}
\usepackage{graphicx}
\usepackage{hyperref}

\usetikzlibrary{automata,positioning}

%
% Basic Document Settings
%

\topmargin=-0.45in
\evensidemargin=0in
\oddsidemargin=0in
\textwidth=6.5in
\textheight=9.0in
\headsep=0.25in

\linespread{1.1}

\pagestyle{fancy}
\lhead{\hmwkAuthorName}
\chead{} % empty
\rhead{\hmwkHeaderTitle}
\cfoot{\thepage}

\renewcommand\headrulewidth{0.4pt}
\renewcommand\footrulewidth{0pt}

\setlength\parindent{0pt}

%
% Homework Details
%

\providecommand{\hmwkTitle}{Homework}
\providecommand{\hmwkDueDate}{TBD}
\providecommand{\hmwkClass}{ICCS206 Discrete Mathematic}
\providecommand{\hmwkClassTime}{Section\ 1}
\providecommand{\hmwkClassInstructor}{Pakawut Jiradilok}
\providecommand{\hmwkAuthorName}{\textbf{Jiraroj Wiruchpongsanon (6781617)}}
\providecommand{\hmwkHeaderTitle}{\hmwkClass:\ \hmwkTitle}

%
% Title Page
%

\title{
    \vspace{2in}
    \textmd{\textbf{\hmwkClass:\ \hmwkTitle}}\\
    \normalsize\vspace{0.1in}\small{Due\ on\ \hmwkDueDate}\\
    \vspace{0.1in}\large{\textit{\hmwkClassInstructor\ \hmwkClassTime}}
    \vspace{3in}
}

\author{\hmwkAuthorName}
\date{}

\renewcommand{\part}[1]{\textbf{\large Part \Alph{partCounter}}\stepcounter{partCounter}\\}

%
% Helper Commands
%

\newcounter{partCounter}
\newcommand{\solution}{\textbf{\large Solution}}

\newtheorem*{theorem}{Theorem}

\theoremstyle{remark}
\newtheorem*{remark}{Remark}

\begin{document}

\maketitle
\newpage


\section*{Problem 1: Basic Stuff}
\begin{enumerate}
    \item[1.1)] Write down the set of integers that is less than $8$ and greater than $2$.
    \item[1.2)] Write down the members of $\{x \in I \mid x^2 < 10\}$.
    \item[1.3)] Write down the members of $\{x \in I^+ \mid x^2 < 10\} \cup \{x \in I^+ \mid 2 < x < 8\}$.
    \item[1.4)] Write down the members of $\{x \in I^+ \mid x^2 < 10\} \cap \{x \in I^+ \mid 2 < x < 8\}$.
\end{enumerate}

\medskip
\solution

\begin{enumerate}
    \item[1.1)] $\{x \in \mathbb{I} \mid 2 < x < 8\}$ or $\{3, 4, 5, 6, 7\}$ explicitely.
    \item[1.2)] $\{-3, -2, -1, 0, 1, 2, 3\}$
    \item[1.3)] $\{1, 2, 3\} \cup \{3, 4, 5, 6, 7\} = \boxed{\{1, 2, 3, 4, 5, 6, 7\}}$ 
    \item[1.4)] $\{1, 2, 3\} \cap \{3, 4, 5, 6, 7\} = \boxed{\{3\}}$
\end{enumerate}


\vspace{1em}
\hrulefill

\section*{Problem 2: Notation Exercise}
Let
\begin{itemize}
    \item $M$ be the set of all MUIC students,
    \item $P(x)$ be the predicate ``student $x$ takes Discrete Math this term'',
    \item $Q(x)$ be the predicate ``student $x$ takes Data Structure this term'',
    \item $F(x,y)$ be the predicate ``students $x$ and $y$ are friends on Facebook'',
    \item $A = \{x \in M \mid P(x)\}$,
    \item $B = \{x \in M \mid Q(x)\}$.
\end{itemize}
Write the following propositions using $\cap$, $\cup$, $\rightarrow$, $\forall$, $\exists$ and set notation.
\begin{enumerate}
    \item[2.1)] There are some students at MUIC who take Discrete Math this term. (Example: $\exists x \in M$ such that $P(x)$.)
    \item[2.2)] There are some students who take both Discrete Math and Data Structure this term.
    \item[2.3)] There exist students who take Discrete Math but not Data Structure.
    \item[2.4)] Everyone takes Discrete Math.
    \item[2.5)] No one takes Data Structure.
    \item[2.6)] All students who take Discrete Math this term also take Data Structure. (Example: $\forall x \in A,\ Q(x)$ or $\forall x \in M,\ P(x) \rightarrow Q(x)$.)
    \item[2.7)] There exist students who take Data Structure but not Discrete Math.
    \item[2.8)] There exists a student in Discrete Math who is friends with everyone in Data Structure.
    \item[2.9)] Everyone in Data Structure has at least one friend in Discrete Math.
    \item[2.10)] There exists a student in Discrete Math who is friends with no one in Data Structure.
\end{enumerate}

\medskip
\solution

\begin{enumerate}
    \item[2.2)] $\exits x \in M$ such that $P(x)$ and $Q(x)$; or $\exists(A \cap B)$
    \item[2.3)] $\exists x \in M \text{ such that } P(x) \land \neg Q(x)$; or $(A \setminus B) \neq \varnothing$; or $\exists x \in A \cap B^{c}$.
    \item[2.4)] $\forall x \in M,\ P(x)$; or $M = A$
    \item[2.5)] $B = \emptyset$
    \item[2.7)] $\exists x \in M \text{ such that } Q(x) \land \neg P(x)$
    \item[2.8)] $\{\text{facebook friend}\} \subseteq \{\text{friend}\}$ --- since you have to be friend first, to be a facebook friend --- a facebook friend counts as a friend, but a friend might not be facebook friend --- though i have neither an active facebook account nor a friend.
        \[\boxed{\exists x \in A \text{ such that } \forall y \in B,\ F(x, y)}\]
    \item[2.9)] $\forall x \in B \text{ such that } \exists y \in B,\ F(x, y)$
    \item[2.10)] $\exists x \in A \text{ such that } \forall y \in B,\ \neg F(x, y)$
\end{enumerate}

\vspace{1em}
\hrulefill

\section*{Problem 3: Truth Tables}
Write down the truth tables for the following statements.
\begin{enumerate}
    \item[3.1)] $P \Rightarrow (\neg Q \lor P)$
    \item[3.2)] $P \Rightarrow (P \land Q)$
    \item[3.3)] $(P \land R) \lor (Q \land R)$
\end{enumerate}

\medskip
\solution

\begin{enumerate}
 \item[(3.1)]
 \[
  \begin{array}{c|c|c|c|c}
  P & Q & \neg Q & \neg Q \lor P & P \Rightarrow (\neg Q \lor P)\\\hline
  T & T & F & T & T\\
  T & F & T & T & T\\
  F & T & F & F & T\\
  F & F & T & T & T
  \end{array}
  \]

 \item[(3.2)]
  \[
  \begin{array}{c|c|c|c}
  P & Q & P \land Q & P \Rightarrow (P \land Q)\\\hline
  T & T & T & T\\
  T & F & F & F\\
  F & T & F & T\\
  F & F & F & T
  \end{array}
  \]

 \item[(3.3)]
  \[
  \begin{array}{c|c|c|c|c|c}
  P & Q & R & P \land R & Q \land R & (P \land R) \lor (Q \land R)\\\hline
  T & T & T & T & T & T\\
  T & T & F & F & F & F\\
  T & F & T & T & F & T\\
  T & F & F & F & F & F\\
  F & T & T & F & T & T\\
  F & T & F & F & F & F\\
  F & F & T & F & F & F\\
  F & F & F & F & F & F
  \end{array}
  \]
\end{enumerate}

\vspace{1em}
\hrulefill

\section*{Problem 4: Fun with Quantifiers}
\begin{enumerate}
    \item[4.1)] Are the following two propositions equivalent?
    \begin{align*}
        &A)~\forall x \in I,\ \exists y \in I \text{ such that } x + y = 23, \\
        &B)~\exists y \in I \text{ such that } x + y = 23,\ \forall x \in I.
    \end{align*}
    \item[4.2)] Let $X$ be the set of all boys in the class and $Y$ the set of all girls. Let $P(a,b)$ be true only if $a$ secretly likes $b$. Translate the following statements into plain English and determine whether they are true. Explain, prove, or disprove each one.
    \begin{enumerate}
        \item[(A)] $\forall x \in X,\ \exists y \in Y \text{ such that } P(x,y) \Rightarrow \exists y \in Y \text{ such that } P(x,y)\ \forall x \in X$
        \item[(B)] $\exists y \in Y \text{ such that } P(x,y)\ \forall x \in X \Rightarrow \forall x \in X,\ \exists y \in Y \text{ such that } P(x,y)$
    \end{enumerate}
\end{enumerate}

\end{document}
